\chapter{Superficies en el espacio}

\begin{definicion}
    Una \textbf{superficie} en $\bb{R}^3$ es un subconjunto $S\subset \bb{R}^3$ no vacío tal que $\forall p \in S$ existe\footnote{cada uno de estos elementos deberían tener subíndice $p$ porque dependen del punto pero por simplificar la notación no lo pondremos}
    \begin{enumerate}
        \item $U\subset \bb{R}^2$ abierto.
        \item un entorno abierto $V$ de $p$ en $S$.
        \item una aplicación $X:U \to \bb{R}^3$ diferenciable.
    \end{enumerate}
    que cumplen:
    \begin{enumerate}
        \item $X(U)=V$
        \item $X:U\to V$ es un homeomorfismo.
        \item $(dX)_q : \bb{R}^2 \to \bb{R}^3$ inyectiva $\forall q=(u,v) \in U$.
    \end{enumerate}
    Normalmente notaremos
    \Func{X}{U}{V}{(u,v)}{(x(u,v), y(u,v), z(u,v))}
    donde $x,y,z:U\to \bb{R}$ son diferenciables.
\end{definicion}

\begin{observacion}
    Con esta definición no puede haber superficies con autointersecciones (ya que en ese caso $X$ no sería biyectivo)\footnote{Realmente sí existen pero no las vamos a estudiar este año.}.
\end{observacion}

\begin{notacion}
    Hemos escrito $(dX)_q$, lo cual hasta ahora habremos notado como $(DX)(q,v)$ que es la derivada direccional en la dirección del vector $v$ del punto $q$. En este caso no especificamos la dirección $v$ porque se considera en todas las direcciones (sticker de shrek).
\end{notacion}

\begin{observacion}
    Otra condición equivalente a que $(dX)_q$ sea inyectiva para todo $q\in U$ sería que para todo $q\in U$ se tenga que
    \begin{gather*}
        \frac{\partial X}{\partial u}(q),\ \ \ \frac{\partial X}{\partial v}(q)\ \ \ \text{ son linealmente independientes}.
    \end{gather*}
    Lo cual además sería equivalente a decir que 
    \begin{gather*}
        \left(\frac{\partial X}{\partial u}(q) \wedge \frac{\partial X}{\partial v}(q)\right) \neq 0\ \ \ \ \forall q\in U
    \end{gather*}
\end{observacion}

\begin{ejemplo}\
    \begin{enumerate}
        \item Cualquier plano $P\subset \bb{R}^3$ se puede escribir como
        \begin{gather*}
            P = \{(x,y,z)\in \bb{R}^3 : ax+by+cz+d=0\}\ \ \ (a,b,c)\neq (0,0,0),\ \ d\in \bb{R}
        \end{gather*}
        Supongamos que $c\neq 0$, entonces 
        \begin{gather*}
            P=\{(x,y,z)\in \bb{R}^3 : z=Ax+By+C\}\ \ \ A,B,C \in \bb{R}
        \end{gather*}
        Veamos que $P$ es una superficie. Para ello, para cada $p\in P$ tenemos que encontrar $U\subset \bb{R}^2$, $V\subset \bb{R}^3$ y $X:U\to V$ que cumplan las condiciones de la definición. Sea $p\in P$ elegimos
        \begin{gather*}
            U=\bb{R}^2, \ \ \ \ V=P
        \end{gather*}
        \Func{X}{U}{\bb{R}^3}{(u,v)}{(u,v,Au+Bv+C)}
        que vemos que no dependen del punto $p$ (ya que estamos en un caso muy sencillo). Además, claramente $X$ es diferenciable. Veamos que se verifican las condiciones para ser una superficie con estos elementos seleccionados:
        \begin{enumerate}
            \item $X(U)=X(\bb{R}^2) = \{(x,y,z)\in \bb{R}^3 : (x,y)\in U, z=Ax+By+C\} = P = V$.
            
            \item Veamos que $X: U=\bb{R}^2 \to V=P$ es un homeomorfismo. Es claramente continua y sobreyectiva, nos falta ver que es inyectiva, lo cual se tiene ya mirando las dos primeras componentes. Podemos considerar su inversa $\pi: P \to \bb{R}^2$ dada por 
            \begin{gather*}
                \pi(x,y,z) = (x,y) \ \ \ \ \forall (x,y,z)\in P
            \end{gather*}
            que es claramente continua. Confirmamos que $X$ es un homeomorfismo.

            \item Se tiene para todo $(u,v)\in U$:
            \begin{gather*}
                X_u(u,v) = (1,0,A)\\
                X_v(u,v) = (0,1,B)
            \end{gather*}
            y es fácil ver que son linealmente independientes (mirando las 2 primeras componentes de cada derivada parcial).
        \end{enumerate}
        \item Veamos que los grafos de funciones son superficies. Para ello consideramos $O\subset \bb{R}^2$, $f:O\to \bb{R}$ diferenciable y tendremos
        \begin{gather*}
            Grafo(f) = \{(x,y,z)\in \bb{R}^3:(x,y)\in O,\ \ z=f(x,y)\}
        \end{gather*}
        Tomamos $p\in Grafo(f)$ y elegimos
        \begin{gather*}
            U=O,\ \ \ \ V=Grafo(f)
        \end{gather*}
        \Func{X}{U}{\bb{R}^3}{(u,v)}{(u,v,f(u,v))}
        donde $X$ es diferenciable por serlo la identidad y $f$. Veamos que cumple las condiciones para ser una superficie:
        \begin{enumerate}
            \item $X(O) = \{(u,v,z)\in \bb{R}^3 : (u,v)\in O,\ \ z=f(u,v)\} = Grafo(f) = V$.
            \item En topología vimos que todo grafo es homeomorfo a su dominio (con la proyección como inversa) por lo que $X$ es un homeomorfismo.
            \item $X_u(u,v) = (1,0,f_u)(u,v)$ y $X_v(u,v)=(0,1,f_v)(u,v)$ que con las dos primeras componentes ya vemos que son linealmente independientes.
        \end{enumerate}
        En este caso vemos de nuevo que estos elementos no dependen del punto $p$ escogido.

        \item Consideramos el siguiente problema:
        \begin{gather*}
            \begin{array}{l}
                X(u,v) = (x,y,z)\\
                X(q) = p
            \end{array}
        \end{gather*}
        \begin{itemize}  
            \item A $(u,v)$ lo llamaremos \textbf{coordenadas} de $p$.
            \item A $X$ lo llamaremos \textbf{parametrización} (local).
            \item A $V$ lo llamaremos \textbf{entorno parametrizado}.
            \item A $U$ lo llamaremos \textbf{dominio de la parametrización}.
        \end{itemize}

        \item Los abiertos de una superficie son una superficie
        
        \item Sea $S\subset \bb{R}^3$ no vacío y supongamos que $S_i\subset S$ es una superficie abierta para cada $i\in I$ tal que
        \begin{gather*}
            S=\bigcup\limits_{i\in I}S_i
        \end{gather*} 
        Entonces $S$ es una superficie.\\

        Para verlo tomamos $p\in S$ por lo que existirá un $i_0\in I$ tal que $p\in S_{i_0}$ superficie. Por tanto tenemos que 
        \begin{gather*}
            \exists U_0\subset \bb{R}^2 \text{ abierto },\ \ \exists V_0 \text{ entorno abierto de $p$ en $S_{i_0}$},\ \ \exists X:U_0\to \bb{R}^3\text{ diferenciable}
        \end{gather*}
        cumpliendo las 3 propiedades de una superficie. Elegimos $V=V_0\subset S$ que es un entorno abierto de $p$ en $S_{i_0}\subset$, que es abierto en $S$. Tomamos $U=U_0$ que es abierto en $\bb{R}$ y tomamos $X_0=X:U=U_0 \to \bb{R}^3$ y se deja como ejercicio ver que se verifican las propiedades de una superficie.

        \item Sea $S\subset \bb{R}^3$ una superficie, $O,O'\subset \bb{R}^3$ abiertos con $S\subset O$ y consideramos $\phi:O\to O'$ un difeomorfismo. Entonces $S'=\phi(S)$ es otra superficie en $\bb{R}^3$.\\
        
        Tomamos $p'\in S'$ y tenemos que $\exists!p\in S$ tal que $\phi(p)=p'$. Como $S$ es superficie, asociados a $p$ existirán $U,V$ y $X$ verificando las propiedades de una superficie. Tomaremos $U'=U$, $V'=\phi(V)$ y $X'=\phi \circ X$. Veamos que cumplen lo que buscamos:
        \begin{enumerate}
            \item $X'(U)=(\phi \circ X)(U)=\phi(X(U)) = \phi(V)=V'$.
            \item Veamos que $X':U'\to V'$. Para ello estamos en la siguiente situación
            \begin{gather*}
                \xymatrix{
                    \pi \circ X: U' \ar[d]  \ar[r] & V'\\
                    V \ar[ur]
                }
            \end{gather*}

            \item Sea $q'\in U'=U$, entonces tenemos que 
            \begin{gather*}
                (dX')_q = d(\phi \circ X)_{q'}
            \end{gather*}
        \end{enumerate}
    \end{enumerate}
\end{ejemplo}

